\chapter{Panic Attacks}

\section*{Definition}
Sudden episodes of intense fear or discomfort that peak within minutes, accompanied by physical symptoms such as palpitations, sweating, trembling, shortness of breath, chest pain, nausea, dizziness, and a sense of impending doom.

\section*{What Happens in Your Body}
Panic attacks involve acute activation of the amygdala and its projections to the hypothalamus, periaqueductal grey, and locus coeruleus, triggering the sympathetic nervous system (fight-or-flight response). Oestrogen modulates the serotonergic and GABAergic systems that inhibit amygdala reactivity; declining oestrogen reduces inhibitory tone, lowering the panic threshold. Reduced progesterone (with its anxiolytic GABA-agonist metabolites) further disinhibits the panic circuit. Respiratory and interoceptive hypersensitivity (catastrophic misinterpretation of bodily sensations) amplifies the panic response.

\section*{Causes}
\begin{itemize}
  \item Panic disorder (recurrent unexpected panic attacks with fear of recurrence)
  \item Perimenopausal and menopausal hormonal fluctuations
  \item Generalised anxiety disorder
  \item Major depressive disorder
  \item Post-traumatic stress disorder
  \item Hyperthyroidism (thyrotoxicosis)
  \item Hypoglycaemia
  \item Mitral valve prolapse
  \item Substance use or withdrawal (caffeine, alcohol, cocaine, cannabis)
  \item Medication side effects (stimulants, thyroid hormone)
\end{itemize}

\section*{When to See a Doctor}
See your doctor if:
\begin{itemize}
  \item Symptoms are severe or getting worse over time
  \item Recurrent panic attacks, panic attacks with suicidal thoughts, new onset after age 40 without clear trigger, or chest pain with shortness of breath (rule out cardiac)
  \item You have other concerning symptoms such as fever, weight loss, or bleeding
\end{itemize}

\section*{Related Symptoms}
\begin{itemize}
  \item Anxiety
  \item Heart palpitations
  \item Shortness of breath
  \item Dizziness
  \item Chest pain
\end{itemize}

\textbf{Important:} If this symptom persists, your doctor will check for serious underlying causes.

\section*{Self-Care / Home Management}
\begin{itemize}
  \item Rest and avoid activities that make it worse.
  \item Maintain adequate hydration and a balanced diet.
  \item Over-the-counter remedies may provide symptomatic relief (consult a pharmacist or doctor).
  \item Monitor symptoms and seek professional advice if they persist or worsen.
  \item Avoid self-medicating with prescription medications without medical guidance.
\end{itemize}

\section*{How Your Doctor Will Evaluate You}
\begin{itemize}
  \item A thorough history and physical examination are the first steps in evaluation.
  \item Laboratory tests (blood work, urinalysis) may be ordered based on clinical suspicion.
  \item Imaging studies (X-ray, ultrasound, CT, MRI) may be indicated when structural pathology is suspected.
  \item Specialist referral is arranged when the diagnosis remains uncertain or requires specific expertise.
\end{itemize}

\section*{Treatment Options}
\begin{itemize}
  \item Treatment targets the underlying cause identified through diagnostic evaluation.
  \item Supportive care and symptom management are provided while awaiting definitive therapy.
  \item Medications, lifestyle modifications, and physical therapies may be prescribed as appropriate.
  \item Follow-up ensures treatment effectiveness and allows timely adjustments to the care plan.
\end{itemize}

\clearpage